%% P22 --- Constrained optimisation and Lagrange multipliers
%%
%% Stub, generated by tools/gen_stubs.py from tools/programs.json.
%% The brief below is the contract for this program: what it must argue, and
%% what it must buy the reader. Writing the program means deleting the
%% \programstub{} block and replacing it with the program -- after which this
%% file is never regenerated. If the brief turns out to be wrong once the
%% mathematics has been worked, change it in the manifest and record the
%% contradiction in CLAUDE.md under *Resolved questions*.

\program{Constrained optimisation and Lagrange multipliers}
\label{prog:P22}

\programstub{%
Argument: at a constrained optimum the gradients of objective and constraint
are parallel; the multiplier is the exchange rate between them. Payoff: **a
Lagrange multiplier is a price** \dash{} how much objective you buy per unit
of constraint relaxed \dash{} and that reading turns the \code{beta} in a
KL-penalised objective from a tuning knob into a quantity with units; the
equivalence between a hard KL constraint and a KL penalty, which is why PPO
and DPO are talking about the same problem; projection onto a set as the
other way to enforce a constraint; KKT conditions stated for recognition,
with the inequality case sketched rather than developed.  FORWARD REFERENCE,
declared: the payoff is stated in terms of a KL-penalised objective, and the
Kullback--Leibler divergence is not defined until P30. State the one fact it
needs -- that KL is a non-negative measure of how far one distribution sits
from another, zero only when they agree -- in the Learning outcomes with the
pointer, exactly as P18 does for cross-entropy. Alternatively carry the
payoff with a plain quadratic penalty and have P30 return to it; that is the
author's call, but it may not be left undeclared.

\medskip
\textbf{Planned length:} 50 frames.
}
